\usepackage{amsmath}
\usepackage{amssymb}
\usepackage{amsthm}
\usepackage{amsfonts}
\usepackage{epic,eepic}
\usepackage[binary-units]{siunitx}
\usepackage[utf8]{inputenc}
\usepackage{stackengine}
\usepackage{graphicx}
\usepackage{hyperref}
\usepackage[capitalise]{cleveref}
\usepackage[]{todonotes}
\usepackage{bbm}
\usepackage{tikz}
\usepackage{circuitikz}
\usepackage{braket}
\usepackage{physics}
\usepackage{fontawesome}
\usepackage{siunitx}
\usetikzlibrary{%
        shapes.geometric,calc,intersections,%
        decorations.markings,decorations.text,decorations,%
        patterns,decorations.pathmorphing,fpu,%
                decorations.pathreplacing,%
        positioning,automata,shapes.multipart,circuits,%
        circuits.ee, circuits.ee.IEC, shapes.gates.ee,%
                arrows.meta,arrows,math,matrix,
		fadings
}

%%%%%%%%%%%%%%%%%%%%%%%
\newcommand{\normord}[1]{\vcentcolon\mathrel{#1}\vcentcolon}
\providecommand{\vcentcolon}{\mathrel{\mathop{:}}}
\newcommand{\mytodo}[1]%
{{\todo[inline,backgroundcolor=blue!10!white]{#1}
}}
\newcommand{\me}{\mathrm{e}}
\newcommand{\mi}{\mathrm{i}}
\newcommand{\md}{\mathrm{d}}
%\DeclareMathOperator{\trace}{tr}
\DeclareMathOperator{\heaviside}{H}
\DeclareMathOperator\Ai{Ai}
\DeclareMathOperator{\sinc}{sinc}
\DeclareMathOperator{\Ein}{Ein}
\DeclareMathOperator{\sign}{sign}

\theoremstyle{definition}
\newtheorem{theorem}{Theorem}[section]
\newtheorem{lemma}{Lemma}[subsection]
\newtheorem{proposition}{Proposition}[section]
\newtheorem{properties}{Properties}[section]
\newtheorem{definition}{Definition}[subsection]
\newtheorem*{remark}{Remark}
\newtheorem*{example}{Example}


%%%%%%%%%%%%%%%%%%%%%%%
%\usepackage[disable]{todonotes}   % uncomment to hide all the notes
\usepackage{todonotes}
\newcommand{\notePDG}[1]{\todo[color=red!40, author=\textbf{Pascal}, inline]{#1}}
\newcommand{\noteBR}[1]{\todo[color=red!40, author=\textbf{Benjamin}, inline]{#1}}
\newcommand{\noteAF}[1]{\todo[color=green!40, author=\textbf{Alex}, inline]{#1}}
%%%%%%%%%%%%%%%%%%%%%%%%
\newenvironment{changemargin}[2]{\begin{list}{}{%
\setlength{\topsep}{0pt}%
\setlength{\leftmargin}{0pt}%
\setlength{\rightmargin}{0pt}%
\setlength{\listparindent}{\parindent}%
\setlength{\itemindent}{\parindent}%
\setlength{\parsep}{0pt plus 1pt}%
\addtolength{\leftmargin}{#1}%
\addtolength{\rightmargin}{#2}%
}\item }{\end{list}}

